\chapter{Operators}
\label{ch:operators}

Chapter~\ref{ch:language} introduced the operator as a grammatical category: a command that waits for a motion or text object before it does anything. This chapter examines the operators themselves, the actual transformations Vim performs once that target is supplied. Where Part II was concerned with the shape of the sentence, this chapter and the two that follow it are concerned with what the verb can mean.

\section{Deletion}

\texttt{d} removes the text spanned by its target and places it in a register (the unnamed register by default, or a named one if specified, per Chapter~10). \texttt{dw} deletes to the beginning of the next word; \texttt{d\$} (or its shorthand \texttt{D}) deletes to the end of the line; \texttt{dd} deletes the current line as a unit, per the doubling convention established in Chapter~\ref{ch:language}. Deletion in Vim is never destructive in the sense of discarding the removed text outright: because the deleted span is always written to a register first, undoing a deletion is rarely even necessary, since the same text can usually be retrieved with a put command instead, a point developed further once registers themselves are covered in Part IV.

\section{Change}

\texttt{c} behaves exactly like \texttt{d} with one addition: after removing the spanned text, it leaves the editor in Insert mode at the point of removal, ready for replacement text to be typed. \texttt{cw} changes to the end of the current word; \texttt{ciw} changes the inner word text object regardless of cursor position within it; \texttt{cc} changes the entire current line, deleting its content but preserving its indentation, which distinguishes \texttt{cc} from a plain \texttt{dd} followed by \texttt{O} in a way that is easy to overlook until it is needed.

The relationship between \texttt{c} and \texttt{d} is worth stating explicitly, because it is a small but genuine exception to the otherwise uniform operator grammar: \texttt{c} is not simply "delete, then separately enter Insert mode," which would be two commands composed by the user, but a single operator whose defined behavior includes that mode transition as part of what the operator itself does. This matters practically because it means \texttt{c} cannot be meaningfully separated into a delete-operator-plus-insert-command pair for the purposes of the period-repeat command from Chapter~\ref{ch:language}: pressing \texttt{.} after \texttt{ciw} followed by typed replacement text repeats the entire compound action, deletion and the subsequent insertion together, against a new target.

\section{Yanking}

\texttt{y} copies the spanned text into a register without removing it from the buffer, the direct analogue of \texttt{d} for the reader who wants to duplicate text rather than relocate it. \texttt{yw} yanks to the beginning of the next word; \texttt{yy} yanks the current line. Because deletion already writes to a register as a side effect, a common point of confusion for newcomers is not understanding why an explicit yank command is needed at all; the answer is that deletion's register write is incidental to removing the text, while yank's entire purpose is that write, with no accompanying removal, and the two operators exist separately because "copy without altering the buffer" and "remove, incidentally retaining a copy" are genuinely different operations that happen to share an underlying storage mechanism.

\section{Case Conversion}

\texttt{gU} uppercases the spanned text and \texttt{gu} lowercases it; \texttt{g\textasciitilde} toggles the case of each character in the span independently. \texttt{gUiw} uppercases the inner word under the cursor; \texttt{guaw} lowercases a word and its surrounding whitespace region (the whitespace itself being unaffected by a case operation, since it contains no letters, but included in the span for consistency with how \texttt{aw} is defined). These operators are a clean illustration of Section~\ref{sec:multiplication}'s multiplicative claim from Chapter~\ref{ch:language}: nothing about case conversion required a special case-conversion grammar distinct from deletion or yanking; it required only a new operator, immediately usable against every motion and text object already known.

\section{Indentation}

\texttt{>} increases indentation by one shift width (configured with the \texttt{shiftwidth} option introduced in Chapter~\ref{ch:installing}) across the spanned lines, and \texttt{<} decreases it. Because indentation is inherently line-oriented, these two operators are conventionally doubled (\texttt{>>}, \texttt{<<}) to affect the current line, and prefixed with a count to affect several lines at once (\texttt{3>>} shifts the next three lines), rather than commonly paired with the finer-grained motions used for deletion or yanking; shifting half a line by one indent level is rarely a meaningful operation, so the indentation operators are, in practice, used almost exclusively in their doubled or Visual-mode forms.

\texttt{=} is a related but distinct operator: rather than shifting by a fixed number of shift widths, it invokes Vim's indentation logic (either its built-in, filetype-aware rules or an external program specified by the \texttt{equalprg} option) to compute what the correct indentation of the spanned lines should be, and applies that instead. \texttt{=i\{} auto-indents the contents of the current code block according to the file's own indentation rules, a genuinely different operation from \texttt{>i\{}, which would shift the block by one fixed amount regardless of what the surrounding code's indentation logic considers correct.

\section{Joining}

\texttt{J} joins the next line onto the end of the current one, inserting a single space at the join point (unless the current line already ends in whitespace, or the next line begins with a closing parenthesis, in which case Vim's built-in joining logic omits the extra space to avoid producing an obviously wrong result). \texttt{gJ} performs the same join without inserting any space, a literal concatenation. A count prefixed to either form joins that many lines at once: \texttt{3J} joins the current line with the following two. Unlike most of the operators in this chapter, \texttt{J} does not take a motion or text object argument; it operates on a fixed unit (the following line, or, with a count, the following several lines) rather than on an arbitrary span, which places it slightly outside the general operator-target grammar even though it shares the operator category by convention.

\section{Filtering Through the Shell}

\texttt{!} is an operator in the formal sense used throughout this chapter, but one whose target is not merely a motion or text object; it sends the spanned text out of Vim entirely, to an external shell command, and replaces it with that command's output. \texttt{!\}sort} takes the current paragraph, pipes it through the \texttt{sort} command, and replaces the paragraph with the sorted result. This operator is the earliest concrete point of contact in this book between the operator grammar covered so far and the wider Unix integration that Part VII treats in depth; it is included in this chapter, rather than deferred entirely, because it obeys the same operator-target grammar as every other transformation here, differing only in where the actual transformation is computed.

\section{Formatting}

\texttt{gq} reformats the spanned lines to wrap at the width specified by the \texttt{textwidth} option, and \texttt{gw} does the same while preserving the cursor's position rather than moving it to the end of the formatted text. \texttt{gqip} reformats the current paragraph, a genuinely common operation when editing prose (covered further in Chapter~20's treatment of writing workflows) where a line has grown too long or a paragraph has been edited into ragged, inconsistent line lengths that no longer wrap cleanly.

\section{Operators as a Closed but Extensible Class}

The operators catalogued in this chapter, deletion, change, yank, case conversion, indentation, joining, shell filtering, and formatting, are not an arbitrary list; they are close to the complete set of built-in, general-purpose transformations Vim applies to a span of text identified by the grammar of Part II. New operators can be defined through Vimscript (Part VI), in exactly the same way new text objects can be, but this book's plugin-free scope means the operators above represent, practically speaking, the full built-in vocabulary a reader needs to combine with every motion and text object already covered. What Chapter~\ref{ch:language} promised as an economical, multiplicative grammar should, by the end of this chapter, feel concretely realized: roughly eight operators, each multiplying against dozens of motions and text objects, already account for several hundred distinct, meaningful commands, the overwhelming majority of which have not been individually demonstrated anywhere in this book and do not need to be.
